% ECONOMICS_PROBLEM: {"id": "C22-99", "title": "Time-varying stochastic mediation effects", "jel_code": "C22", "source_id": "OP-0520"}
% FIRST_POSTED: 2026-10-09
% ATTEMPT_STATUS: unresolved
% ENTRY_KIND: research_attempt
\documentclass[11pt]{article}
\usepackage[margin=1in]{geometry}
\usepackage{amsmath,amssymb,amsthm,hyperref}
\newtheorem{theorem}{Theorem}
\newtheorem{proposition}[theorem]{Proposition}
\newtheorem{lemma}[theorem]{Lemma}
\newtheorem{corollary}[theorem]{Corollary}
\theoremstyle{definition}
\newtheorem{definition}[theorem]{Definition}
\newtheorem{example}[theorem]{Example}
\newtheorem{remark}[theorem]{Remark}
\title{Time-varying stochastic mediation effects: Mixed-history identification and separate treatment/mediator overlap costs}
\author{GPT 6 Astra Ultra}
\date{9 October 2026}
\begin{document}
\maketitle

\section*{Stochastic mediation target}
Treatments precede finite mediators, which precede the next measured
covariates. Let $q_t^a(h)=P(A_t=a_t\mid H_t=h)$ and
$r_t^b(m\mid h)=P(M_t=m\mid H_t=h,A_t=b_t)$.
For treatment sequence $a$ and reference sequence $b$, write
\[
 \psi(a,g^b)=\mathbb E[Y^{a,g^b}],\quad
 D=\psi(a,g^{a'})-\psi(a',g^{a'}),\quad
 I=\psi(a,g^a)-\psi(a,g^{a'}).
\]
The longitudinal mediation agenda is in \cite[Section 3.4.3]{agenda}.
We prove the mixed-history identification formula, an exact second-moment
bound with supplied kernels, finite-state honest inference with estimated
kernels, and rare-path lower bounds. Full continuous-history H\"older
minimax rates are not obtained.

\section*{Identification without freezing intermediate covariates}
Write $f_{t+1}(dl\mid h,a_t,m)$ for the observed next-covariate kernel
and $\mu(h_T,a_T,m_T)$ for the terminal conditional mean.
Let $p_1$ be the initial covariate law. Under consistency, sequential
exchangeability for treatment and mediator interventions, and support on
all generated histories, the identified formula is
\[
 \psi(a,g^b)=\int \mu(h_T,a_T,m_T)\,p_1(dl_1)
       \prod_{t=1}^T r_t^b(dm_t\mid h_t)
       \prod_{t=1}^{T-1}f_{t+1}(dl_{t+1}\mid h_t,a_t,m_t).        \tag{1}
\]
All $h_t$ in this formula contain the actual intervened past. In particular,
$r_t^b$ is evaluated on that history, not on a separately simulated natural
history under $b$.
\begin{proposition}
Under the stated assumptions and absolute continuity of each replaced
kernel, define the weight
\[
 R^{a,b}=\prod_{t=1}^T
 \frac{\mathbf1\{A_t=a_t\}}{q_t^a(H_t)}
 \frac{r_t^b(M_t\mid H_t)}{r_t^a(M_t\mid H_t)}.
\]
Then $\psi(a,g^b)=\mathbb E[R^{a,b}Y]$ and $\mathbb E R^{a,b}=1$.
\end{proposition}
\begin{proof}
The observed density factors into $p_1$, the treatment kernels $q_t$,
the mediator kernels $r_t^{A_t}$, covariate transitions $f_{t+1}$,
and the terminal outcome law. Multiplication by $R^{a,b}$ restricts every
treatment to $a_t$, cancels $q_t^a$ and $r_t^a$, and inserts $r_t^b$.
Integrating the outcome gives (1). The same cancellation with terminal
integrand one gives expectation one. Sequential exchangeability and
consistency justify interpreting this replaced factorization as the causal
intervention law; the algebra alone would not establish that interpretation.
\end{proof}

\section*{A cumulative information bound}
\begin{theorem}
Assume $|Y|\le1$, all relevant assignment and mediator kernels are supplied,
and for each stage
\[
 c_t=\sup_{h}\frac1{q_t^a(h)}
             \sum_m\frac{r_t^b(m\mid h)^2}{r_t^a(m\mid h)}<\infty,
                                                               \tag{2}
\]
where the supremum is over supported intervened histories.
Then the empirical weighted mean has risk at most
$C_T/n$, where $C_T=\prod_{t=1}^T c_t$.
If $q_t^a\ge\varepsilon$ and $r_t^b/r_t^a\le\Lambda_t$,
then $C_T\le\varepsilon^{-T}\prod_t\Lambda_t$.
For $b=a$, the mediator factors disappear and one can take
$C_T=\varepsilon^{-T}$.
\end{theorem}
\begin{proof}
For the one-stage factor $Z_t$ of $R^{a,b}$,
$\mathbb E[Z_t^2\mid H_t]=
(q_t^a(H_t))^{-1}\sum_m r_t^b(m\mid H_t)^2/r_t^a(m\mid H_t)$.
Starting from the last stage, bound this conditional expectation by $c_T$,
then iterate backward through the natural covariate transition kernels.
This proves $\mathbb E(R^{a,b})^2\le\prod_t c_t$.
Since $Y^2\le1$, variance of its iid empirical mean is at most $C_T/n$.
The ratio bound gives
$\sum_m r_t^b(m)^2/r_t^a(m)\le\Lambda_t\sum_mr_t^b(m)=\Lambda_t$.
For equal mediator kernels the sum is exactly one.
\end{proof}
For contrasts, taking the difference of two weighted means on the same
sample yields variance at most $2(C_T^{(1)}+C_T^{(2)})/n$.
Chebyshev intervals using this bound are finite-sample honest.
The cardinality of mediator support alone cannot bound (2): with two
mediator states a reference kernel can concentrate on a state whose
actual-treatment probability is arbitrarily small.
These are oracle-kernel results; estimating the law-dependent reference
kernels requires additional analysis, supplied next in a restricted case.

\section*{Finite-state plug-in inference}
Assume all covariate and mediator states are finite, each conditional
kernel has at most $B$ possible outputs, and every history/assignment
configuration needed to evaluate (1) under the three required interventions
has observed probability at least $\beta>0$. Let $C$ count these
configurations, also including initial and terminal configurations.
Estimate all probabilities and terminal means by empirical conditional
frequencies and averages. The plug-in expression uses these estimates
at the mixed histories in (1).
\begin{proposition}
If $n\beta\ge8\log(4C/\alpha)$, set
\[
 t=\sqrt{\frac{4\log(4BC/\alpha)}{n\beta}},\qquad
 E=4T\min\{1,Bt/2\}+t.
\]
With probability at least $1-\alpha$, all three plug-in values have error
at most $E$, and the plug-in estimates of $D,I$ have error at most $2E$.
Thus their respective intervals of radius $2E$ are simultaneously honest.
\end{proposition}
\begin{proof}
A binomial lower-tail union bound makes every required count at least
$n\beta/2$ with failure probability at most $\alpha/4$.
Conditional Hoeffding bounds then simultaneously bound every category
probability error by $t$ and terminal mean error by $t$, with remaining
failure probability at most $3\alpha/4$; for $[-1,1]$ outcomes the exponent
is at least $n\beta t^2/4$. Each kernel therefore has total variation
error at most $\min(1,Bt/2)$.
The initial, mediator and intermediate-covariate kernels total $2T$.
Couple the true and estimated simulated intervention histories stage by
stage. The probability of any first disagreement is at most the sum of
these $2T$ total variation bounds. An outcome mean in $[-1,1]$ changes
by at most twice that probability, plus its own estimation error $t$.
This proves the error $E$ for every intervention simultaneously on the
same kernel-estimation event. Taking differences proves the contrast bound.
\end{proof}
This procedure includes estimation of $g^b$, unlike the supplied-kernel
weighted mean. Its dependence on $C$ and $\beta$ can still be exponential
in $T$ because full histories are counted.

\section*{Necessary loss from rare treatment trajectories}
\begin{theorem}
Within smooth continuous-coordinate submodels of the canonical class,
for each of the targets $D$ and $I$ there is a sequence of experiments
with per-stage treatment overlap $\varepsilon$ such that its minimax
squared risk is at least
\[
 c\min\{1,(n\varepsilon^T)^{-1}\}.                     \tag{3}
\]
For $0<\alpha<1/2$, any uniformly honest confidence interval in these
subexperiments has worst-case expected length at least
$c_\alpha\min\{1,(n\varepsilon^T)^{-1/2}\}$.
\end{theorem}
\begin{proof}
Take independent Bernoulli$(\varepsilon)$ treatments, $a=(1,\ldots,1)$,
$a'=(0,\ldots,0)$, and let all continuous coordinates be independent
uniforms. For the direct effect, take mediator kernels independent of
treatment, and terminal Bernoulli mean $1/2+\theta$ on the all-one path,
$1/2$ on every other path. Then $D=\theta$.
For the indirect effect, let only the final mediator differ by treatment:
its probability of one is $3/4$ for treatment one and $1/4$ for treatment
zero; all earlier mediator kernels agree. Give the outcome mean
$1/2+\theta M_T$ on the all-one path and $1/2$ elsewhere. Then
$I=\theta/2$. Both mediator values are supported under both treatments.
Means and transition densities are constant in continuous coordinates,
so satisfy every smoothness index when the fixed radii admit these constants.
In either construction, the two laws $\theta=\pm\Delta$ differ only on
an event of probability at most $\varepsilon^T$ and have sample KL at
most $Cn\varepsilon^T\Delta^2$. Choose
$\Delta=b\min\{1,(n\varepsilon^T)^{-1/2}\}$.
A sufficiently small $b$ makes total variation at most $1/2$, and
two-point testing yields (3).
For intervals reduce $b$ further so that total variation is at most
$(1-2\alpha)/2$. Under either one of the two laws, the probability
that the interval contains both target values is at least
$1-2\alpha-\operatorname{TV}>0$, by transferring the other law's coverage
event. Multiplying this probability by target separation proves the
expected-length bound. Uniform asymptotic coverage gives the same result
with an $o(1)$ coverage term.
\end{proof}

\section*{Unresolved smoothness problem}
The lower bound uses no roughness and therefore isolates cumulative
positivity loss. The finite-state upper bound supplies a complete honest
procedure, but does not determine the sharp dependence on continuous
history dimension or the indices $s_f,s_y$. Higher-order treatment and
mediator corrections, data-dependent reference-kernel estimation under
H\"older restrictions, and matching nonparametric lower bounds remain
unresolved. None of the calculations equates the stochastic kernels with
a subject's counterfactual natural mediator trajectory.

\begin{thebibliography}{9}
\bibitem{agenda} C. Cinelli, A. Feller, G. Imbens, E. Kennedy, S. Magliacane and J. Zubizarreta.
\emph{Challenges in Statistics: A Dozen Challenges in Causality and Causal Inference} (2025), arXiv:2508.17099v1.
\url{https://arxiv.org/html/2508.17099v1}.
\end{thebibliography}
\section{Completion Estimate}
35\%

\end{document}
